% Abstract fragment for a venue template; not a compiled manuscript.
\begin{abstract}
Local large language model services face changing compute availability and
communication conditions, but redistributing model layers can incur weight
loading and key-value cache reconstruction costs. We present Shardwise, a
standalone prototype that combines contiguous-layer partitioning, live telemetry,
hysteresis and an explicit transition gate. The controller compares useful-token
capacity over a planning horizon rather than accepting every improved
steady-state placement. We evaluate its measurement and decision mechanisms
through separated endpoint/context cost ablations, transition-cost sensitivity,
and a controlled worker-level queueing pilot. Decision-only experiments show
that removing endpoint or context terms changes selected placements under a
common reference model; these differences are predictions rather than measured
speedups. Across 54 instrumented loopback windows and 66,003 RPCs, a 50 ms service at
four-client concurrency produces a 150.0 ms RPC-minus-compute residual,
of which 99.2\% is measured lock waiting. Thus this signal can
misattribute concurrency-induced delay to communication. Bounded FIFO admission
and direct lock-wait measurements separate this effect while retaining client
waiting in end-to-end metrics. The study provides reproducible ablations and
identifies calibration requirements for adaptive inference. Demonstrating that
repartitioning improves real distributed service, including transition downtime,
remains the next experimental step.
\end{abstract}
